\documentclass[11pt]{article}
\usepackage{preamble}
\renewcommand{\answer}[1]{}

\begin{document}

\ladderheading{AP Statistics \textperiodcentered\ Unit 1 \textperiodcentered\ Topic 1.4 \textperiodcentered\ B: 2026-09-14 \textperiodcentered\ E: 2026-09-16}{Same Counts, Different Message}

\focusbanner{TODAY'S PROBLEM}

Suppose 184 students chose a favorite lunch option. A newsletter printed \textbf{Graph A} with the headline ``\textbf{Pizza wins by a landslide.}'' A council member redrew the data as \textbf{Graph B} and said, ``It is basically a tie.''\par
\begin{center}
\begin{minipage}{0.48\linewidth}
\centering
\begin{tikzpicture}
\begin{axis}[
  width=0.95\linewidth,
  height=1.4in,
  ybar=2pt,
  bar width=10pt,
  ymin=40,
  ymax=55,
  ytick={40,44,48,52},
  symbolic x coords={Pizza,Tacos,Salad,Pasta},
  xtick=data,
  enlarge x limits=0.18,
  title={\textbf{Graph A}},
  xlabel={Favorite lunch option},
  ylabel={Number of students},
  x tick label style={rotate=30,anchor=east,font=\scriptsize},
  tick label style={font=\scriptsize},
  label style={font=\scriptsize},
  title style={font=\small},
  ymajorgrids,
  nodes near coords,
  every node near coord/.append style={font=\scriptsize},
  axis on top
]
\addplot[fill=calloutblue,draw=dayblue] coordinates
  {(Pizza,52) (Tacos,47) (Salad,44) (Pasta,41)};
\end{axis}
\end{tikzpicture}
\end{minipage}\hfill
\begin{minipage}{0.48\linewidth}
\centering
\begin{tikzpicture}
\begin{axis}[
  width=0.95\linewidth,
  height=1.4in,
  ybar=2pt,
  bar width=10pt,
  ymin=0,
  ymax=58,
  ytick={0,10,20,30,40,50},
  symbolic x coords={Pizza,Tacos,Salad,Pasta},
  xtick=data,
  enlarge x limits=0.18,
  title={\textbf{Graph B}},
  xlabel={Favorite lunch option},
  ylabel={Number of students},
  x tick label style={rotate=30,anchor=east,font=\scriptsize},
  tick label style={font=\scriptsize},
  label style={font=\scriptsize},
  title style={font=\small},
  ymajorgrids,
  nodes near coords,
  every node near coord/.append style={font=\scriptsize},
  axis on top
]
\addplot[fill=calloutblue,draw=dayblue] coordinates
  {(Pizza,52) (Tacos,47) (Salad,44) (Pasta,41)};
\end{axis}
\end{tikzpicture}
\end{minipage}
\end{center}
\begin{firsttakebox}{FIRST TAKE --- before the video (5 min)}
Is pizza clearly the favorite? Name the graph that most influenced you.
\workspace{0.9in}
\end{firsttakebox}

\begin{callout}{\IconBook\ GRAPHING ONE CATEGORICAL VARIABLE (keep on the page)}{calloutgreen}
In a bar graph, each bar's height represents a category's \textbf{count} or \textbf{relative frequency};
bars have equal widths and gaps. Bar height is measured from the displayed baseline, so changing the baseline
changes the visible height even when the count does not change.\\
Relative frequency $=$ count $\div$ total. A \textbf{bar graph} encodes categories as aligned lengths.
A \textbf{pie chart} encodes shares of one whole as angles and areas. Match the display to the reader's question.
\end{callout}

\newpage
\textbf{Finish after the video:}\par\smallskip

\dokbadge{1}~\textbf{(a)} Read the Pizza and Pasta counts. How many more students chose Pizza?\par
\workspace{0.7in}

\dokbadge{2}~\textbf{(b)} Convert all four counts to whole-number percents of 184. Describe how close the categories are in context.\par
\workspace{1.4in}

\dokbadge{3}~\textbf{(c)} For ``Is pizza clearly the favorite?'' choose the more honest graph and name the design feature that settles it. Quantify Graph A's visible Pizza-to-Pasta heights, then compare 52 with 41 using percent increase. Finally, choose a zero-baseline bar graph or pie chart for comparing the four categories and defend that choice.\par
\workspace{2.0in}

\begin{sentenceframebox}\raggedright\textbf{Frame:} Graph \blank[0.5in] is more honest because its axis begins at \blank[0.4in], so height represents \blank[1.2in].\end{sentenceframebox}

\begin{sentenceframebox}\raggedright\textbf{Frame:} Graph A looks \blank[0.5in] times as tall, but 52 is only about \blank[0.5in]\% more than 41.\end{sentenceframebox}

\begin{summaryexitbox}{EXIT --- turn this in}
One thing the video changed about my first take: \hrulefill\par
\workspace{0.4in}
\end{summaryexitbox}

\end{document}
